% ==================== 第 2 章 相关技术与理论基础 ====================
\chapter{相关技术与理论基础}
\label{ch:background}

本章综述本文所依赖的核心技术与理论，为后续章节的设计与实现奠定基础。内容涵盖对称密码与认证加密、密钥派生与信封加密、密钥管理系统、可信执行环境以及身份认证与访问控制五个方面，并在每一方面指出本文的具体取舍。

\section{对称密码与认证加密}
\label{sec:aead}

对称密码使用同一密钥进行加密与解密，具有运算高效的特点，适合加密大量数据。然而，仅保证机密性的加密模式（如 CBC）无法抵御密文篡改，历史上由此产生了诸如填充预言（padding oracle）等攻击。认证加密（Authenticated Encryption with Associated Data, AEAD）在加密的同时提供完整性与真实性保证，并允许将一段“关联数据”绑定到密文上参与完整性校验，但该关联数据本身不被加密\cite{rogaway2002aead, rfc5116}。

本文统一采用 AES-256-GCM 作为认证加密算法\cite{fips197, dworkin2007gcm, mcgrew2004gcm}。GCM（Galois/Counter Mode）在计数器模式加密的基础上引入 GHASH 认证，输出密文与一个认证标签（tag）；解密时若标签校验失败则拒绝输出明文，从而同时保证机密性与完整性。使用 GCM 时需特别注意两点：其一，随机数（nonce）在同一密钥下必须唯一，nonce 重用会严重削弱安全性，因此本文在每次加密时均生成新的 96 位随机 nonce；其二，关联数据可用于将密文与其使用上下文绑定，本文借此把数据密文与具体密钥身份绑定（详见第~\ref{ch:design}~章）。选择 GCM 而非 CBC 的根本原因即在于其 AEAD 特性，可一并解决机密性与完整性问题，并规避填充预言一类风险。

\section{密钥派生与信封加密}
\label{sec:envelope}

在实际系统中，直接用单一密钥加密所有数据会带来密钥暴露面过大、轮换困难等问题。信封加密（envelope encryption）通过分层结构缓解这一矛盾：系统持有一把高价值的主密钥（Master Key），用它“包装”（加密）若干数据密钥（Data Encryption Key, DEK），再以 DEK 加密业务数据。被包装后的 DEK 可安全地与密文一同存储，而主密钥仅用于包装/解包，使用频率与暴露面大大降低。这一模式被各大云 KMS 广泛采用\cite{nist80057}。

与系统密钥相关但目标不同的，是用户口令的处理。口令不应被加密存储，而应使用专门的口令哈希函数进行单向保护。本文对用户口令采用 bcrypt\cite{provos1999bcrypt}，它通过加盐与可调代价因子抵抗暴力破解；这与用于派生系统密钥的 PBKDF2\cite{rfc8018} 等密钥派生函数在用途上有本质区别——前者用于验证口令，后者用于生成密钥材料。明确区分“口令哈希”与“密钥派生/包装”是密钥管理工程中的基本原则。

\section{密钥管理系统}
\label{sec:kms}

密钥管理系统在公钥基础设施（PKI）与云安全中承担密钥全生命周期的集中管理职责。一个完整的密钥生命周期通常包括创建（Create）、使用（Use）、轮换（Rotate）、吊销（Revoke）与销毁（Destroy）等阶段\cite{nist80057}。轮换用于限制单把密钥的使用时长与数据量，吊销与销毁用于在密钥可能泄露或不再需要时及时止损；其中轮换需兼顾“新数据用新密钥”与“历史密文仍可解密”这一对看似矛盾的需求，通常通过保留密钥的历史版本来实现。

在产品形态上，KMS 可由硬件安全模块（HSM）提供最强的密钥隔离，也可由云 KMS（如各类公有云密钥服务）以托管方式提供弹性与易用性。本文实现的系统在定位上不同于上述生产级方案：它是一个面向教学与研究的\textbf{轻量级工程原型}，强调架构清晰、可复现与威胁模型的诚实表述，而非高可用、合规认证或硬件级隔离。第~\ref{ch:eval}~章将对本文系统与云 KMS 在功能、威胁模型与适用场景上作定性对比。

\section{可信执行环境}
\label{sec:tee}

可信执行环境（TEE）的基本思想是在不可信的主机环境中，借助硬件机制构造一块受保护的执行区域（通常称为 Enclave），使其中的代码与数据的机密性和完整性不依赖于操作系统或更高权限的软件\cite{sabt2015tee}。Intel SGX 通过受保护的内存区域（Enclave Page Cache）与一系列处理器指令实现隔离执行\cite{costan2016sgx, mckeen2013sgx}；ARM TrustZone 则通过将系统划分为“安全世界”与“普通世界”来隔离敏感计算\cite{pinto2019trustzone}。

TEE 的两项关键能力与本文密切相关。其一是\textbf{远程认证}（Remote Attestation）：Enclave 能够生成一份可被远端验证的、关于其自身代码与配置的度量证明，使验证方确信自己正在与一个预期的、未被篡改的 Enclave 通信；在 SGX 中这一过程涉及 \texttt{EREPORT}/\texttt{QUOTE} 等机制，并由硬件厂商的认证服务背书\cite{anati2013attestation}。其二是\textbf{密封存储}（Sealed Storage）：Enclave 可用与其身份绑定的密钥加密数据后存盘，使得只有相同度量值的 Enclave 才能解封。

\textbf{本文采用软件模拟 TEE。}具体而言，本文以一个独立的操作系统进程承载 Enclave，对隔离执行、远程认证与安全信道等能力进行结构性模拟：以代码度量值的比对模拟身份验证，以挑战-响应与会话密钥协商模拟认证握手，以 AES-GCM 信道模拟安全通信。必须强调，软件模拟与真实硬件 TEE 存在本质差异：模拟 Enclave 仍是普通进程，不具备硬件内存加密与隔离，其远程认证也不具备硬件根信任。这些差异构成本文系统的能力边界，将在第~\ref{ch:threat}~章作出明确声明。

\section{身份认证与访问控制}
\label{sec:authz}

KMS 必须严格控制“谁能对哪些密钥执行何种操作”。本文在身份认证层面采用基于 JWT（JSON Web Token）的 Bearer 令牌机制\cite{rfc7519}：用户登录后获得一个带签名与过期时间的令牌，后续请求通过携带该令牌完成无状态鉴权。令牌签名采用 HMAC-SHA256\cite{rfc2104, fips180}，其密钥由服务端持有。

在授权层面，本文采用基于角色的访问控制（Role-Based Access Control, RBAC），定义管理员（Admin）、密钥管理员（Key Manager）与应用用户（App User）三类角色，按角色限制接口访问权限（详见第~\ref{ch:design}~章的权限矩阵）。此外，审计日志记录关键操作的执行者、动作、目标与结果，为事后追溯与异常检测提供依据，是安全系统不可或缺的组成部分。

\section{本章小结}

本章梳理了本文所依赖的关键技术：以 AES-256-GCM 实现认证加密，以信封加密构建分层密钥结构，以密钥生命周期管理组织 KMS 功能，以 TEE 思想隔离敏感操作，并以 JWT、RBAC 与审计构建访问控制与可追溯能力。其中，本文对 TEE 采用软件模拟这一关键取舍贯穿全文，其能力边界将在下一章的威胁模型中被严格界定。
